% =============================================================================
% 1.5 Recherche d'information et modèles de classement
% =============================================================================
\section{Recherche d'information et modèles de classement}
\label{sec:ch1-recherche}

\subsection{Fondements de la recherche d'information}

La recherche d'information est une discipline fondatrice du traitement
automatisé des documents, dont les origines remontent à la vision du
\anglais{Memex} formulée par Vannevar Bush en 1945 et aux premiers systèmes
documentaires des années 1960. Elle répond à une question unique~: étant donnée
une requête, comment identifier et ordonner les documents les plus pertinents
d'une collection~?

Un système de recherche s'articule en quatre opérations~: l'\btpterm{indexation}
transforme les documents en une représentation interrogeable~;
l'\btpterm{interrogation} transforme la requête dans le même espace~; le
\btpterm{classement} ordonne les candidats par pertinence estimée~;
l'\btpterm{évaluation} mesure la qualité de ce classement au regard d'un
jugement de référence.

\subsubsection{Métriques d'évaluation}

\begin{btptable}{Métriques d'évaluation de la recherche d'information}
\label{tab:metriques-ri}
\begin{tabularx}{\linewidth}{@{}L{2.8cm} Y L{3.1cm}@{}}
\toprule
\rowcolor{btppale}
\btpth{Métrique} & \btpth{Définition} & \btpth{Ce qu'elle mesure} \\
\midrule
Rappel à $K$ & Proportion des documents pertinents présents dans les $K$ premiers résultats. & Couverture~: le système laisse-t-il échapper des sources~? \\
Précision à $K$ & Proportion de documents pertinents parmi les $K$ premiers résultats. & Pureté~: le contexte transmis au modèle est-il exploitable~? \\
Rang réciproque moyen & Moyenne de l'inverse du rang du premier document pertinent. & Qualité de la tête de classement. \\
Gain cumulé normalisé & Gain cumulé escompté, pondéré par la position et normalisé. & Qualité du classement avec gradation de pertinence. \\
\bottomrule
\end{tabularx}
\end{btptable}

Pour notre application, la distinction entre rappel et précision est
opérationnelle. Le rappel gouverne l'exhaustivité~: un passage pertinent non
récupéré ne pourra jamais être cité. La précision gouverne la qualité du
contexte transmis au générateur~: des extraits hors sujet dégradent la
génération et augmentent le coût d'inférence.

\subsection{Modèles lexicaux}

\subsubsection{Le modèle vectoriel et la pondération TF-IDF}

Formalisé par Salton, le modèle vectoriel représente chaque document et chaque
requête comme un vecteur dont les dimensions sont les termes du vocabulaire. La
pertinence est estimée par la similarité cosinus entre les deux vecteurs. La
pondération TF-IDF combine deux intuitions~: un terme fréquent dans un document
y est important (\anglais{term frequency}), mais un terme fréquent dans toute la
collection est peu discriminant (\anglais{inverse document frequency}).

\subsubsection{Le modèle probabiliste BM25}

Développé par Robertson et ses collaborateurs à partir des travaux fondateurs de
Robertson et Sparck Jones (1976), le modèle BM25 s'inscrit dans la théorie
probabiliste de la recherche d'information. Il demeure aujourd'hui l'un des
modèles lexicaux les plus performants et sert de référence dans la plupart des
campagnes d'évaluation.

Le score attribué à un document $D$ pour une requête $Q$ s'écrit~:

\begin{equation}
  \mathrm{score}(D, Q) = \sum_{q_i \in Q}
    \mathrm{IDF}(q_i) \cdot
    \frac{f(q_i, D)\,(k_1 + 1)}
         {f(q_i, D) + k_1 \left(1 - b + b\,\dfrac{|D|}{\mathrm{avgdl}}\right)}
  \label{eq:bm25}
\end{equation}

où $f(q_i, D)$ désigne la fréquence du terme $q_i$ dans $D$, $|D|$ la longueur
du document, $\mathrm{avgdl}$ la longueur moyenne des documents de la
collection, $k_1$ un paramètre de saturation de la fréquence (usuellement entre
$1{,}2$ et $2{,}0$) et $b$ un paramètre de normalisation par la longueur
(usuellement $0{,}75$).

Trois propriétés de l'équation~\eqref{eq:bm25} méritent d'être soulignées, car
elles expliquent sa robustesse~:

\begin{itemize}
  \item \btpterm{Saturation de la fréquence.} Le terme en $f/(f + k_1 \cdots)$
        est borné~: la centième occurrence d'un mot apporte beaucoup moins que
        la deuxième. Le modèle résiste ainsi à la répétition artificielle.
  \item \btpterm{Normalisation souple par la longueur.} Le paramètre $b$ permet
        de moduler la pénalisation des documents longs, au lieu de l'imposer.
  \item \btpterm{Pondération par la rareté.} Le facteur $\mathrm{IDF}$ accorde
        un poids déterminant aux termes rares, ce qui est précisément le
        comportement attendu pour une référence normative ou un numéro d'article.
\end{itemize}

\subsubsection{Limites des approches lexicales}

Malgré leur efficacité, les approches lexicales présentent des limites
structurelles~: la \btpterm{synonymie} n'est pas traitée, « voiture » et
« automobile » restant étrangers l'un à l'autre~; la \btpterm{polysémie} n'est
pas résolue, un même terme recouvrant des sens distincts~; enfin la
\btpterm{sensibilité à la formulation} fait qu'une même question posée
différemment produit des résultats différents.

\subsection{Modèles sémantiques}

L'émergence des représentations vectorielles continues a permis de dépasser ces
limites en capturant la sémantique plutôt que la forme. Les plongements
statiques (Word2Vec, GloVe) ont d'abord établi qu'un mot peut être représenté
par un vecteur dans un espace sémantique. Les plongements contextuels, avec
BERT et ses successeurs, ont ensuite rendu cette représentation dépendante du
contexte d'occurrence, résolvant en grande partie la polysémie.

Pour la recherche documentaire, des modèles de plongement spécialisés ont été
développés, dont le tableau~\ref{tab:modeles-plongement} donne un aperçu.

\begin{btptable}{Familles de modèles de plongement pour la recherche}
\label{tab:modeles-plongement}
\begin{tabularx}{\linewidth}{@{}L{2.6cm} Y L{3.2cm}@{}}
\toprule
\rowcolor{btppale}
\btpth{Famille} & \btpth{Caractéristiques} & \btpth{Intérêt} \\
\midrule
Modèles à représentation mixte & Produisent conjointement un vecteur dense et une représentation creuse. & Recherche hybride sans double indexation. \\
Modèles multilingues alignés & Espace sémantique partagé entre langues. & Corpus mêlant français et anglais. \\
Modèles à instruction & La requête et le document sont encodés avec des préfixes distincts. & Meilleure asymétrie requête-document. \\
\bottomrule
\end{tabularx}
\end{btptable}

\subsection{Recherche hybride et fusion de classements}

\subsubsection{Motivation}

Les approches lexicales et sémantiques sont complémentaires plutôt que
concurrentes. La première excelle sur les termes exacts, les identifiants et les
références normatives~; la seconde sur les concepts, les synonymes et les
reformulations. Le tableau~\ref{tab:complementarite} illustre cette
complémentarité sur des requêtes typiques du domaine.

\begin{btptable}{Complémentarité des approches sur des requêtes d'appel d'offres}
\label{tab:complementarite}
\begin{tabularx}{\linewidth}{@{}Y C{2.6cm} C{2.6cm}@{}}
\toprule
\rowcolor{btppale}
\btpth{Type de requête} & \btpth{Lexical (BM25)} & \btpth{Sémantique (dense)} \\
\midrule
« ISO 9001 »                        & \cotefort  & \cotemoyen \\
« Article 5.2 du CCAP »             & \cotefort  & \cotefaible \\
« pénalités de retard »             & \cotemoyen & \cotefort \\
« garantie décennale »              & \cotemoyen & \cotefort \\
« expérience en assainissement »    & \cotefaible & \cotefort \\
\bottomrule
\end{tabularx}
\end{btptable}

\subsubsection{La fusion par rang réciproque}

Combiner les résultats de plusieurs systèmes est un problème central de la
recherche d'information, étudié depuis les travaux de Fox sur la fusion de
données et ceux de Lee sur la combinaison d'évidences. Parmi les méthodes
proposées, la \btpterm{fusion par rang réciproque}, introduite par Cormack,
Clarke et Buettcher en 2009, se distingue par sa simplicité et son efficacité
empirique.

Le score fusionné d'un document $D$ s'écrit~:

\begin{equation}
  \mathrm{score}_{\mathrm{RRF}}(D) =
  \sum_{i \in \mathcal{S}} \frac{1}{k + r_i(D)}
  \label{eq:rrf}
\end{equation}

où $\mathcal{S}$ désigne l'ensemble des systèmes fusionnés, $r_i(D)$ le rang de
$D$ dans le classement du système $i$, et $k$ une constante d'amortissement
fixée usuellement à $60$.

\begin{btpbox}[Pourquoi la fusion par rang réciproque fonctionne]
L'équation~\eqref{eq:rrf} n'utilise que les \btpterm{rangs}, jamais les scores.
Cette propriété résout élégamment le problème d'échelle~: un score BM25, non
borné, et une similarité cosinus, comprise entre $-1$ et $1$, ne sont pas
comparables et leur normalisation est arbitraire. La formule est en outre
robuste au bruit, car l'amortissement par $k$ écrase les écarts entre rangs
élevés~: une erreur ponctuelle de classement a peu d'effet sur le résultat
fusionné.
\end{btpbox}

Le tableau~\ref{tab:exemple-rrf} illustre le calcul sur trois documents.

\begin{btptable}{Exemple de calcul de fusion par rang réciproque ($k = 60$)}
\label{tab:exemple-rrf}
\begin{tabularx}{\linewidth}{@{}L{2.2cm} C{2.2cm} C{2.6cm} Y@{}}
\toprule
\rowcolor{btppale}
\btpth{Document} & \btpth{Rang BM25} & \btpth{Rang sémantique} & \btpth{Score fusionné} \\
\midrule
Document A & 1 & 3 & $1/61 + 1/63 = 0{,}03228$ \\
Document B & 2 & 2 & $1/62 + 1/62 = 0{,}03226$ \\
Document C & 5 & 1 & $1/65 + 1/61 = 0{,}03178$ \\
\bottomrule
\end{tabularx}
\end{btptable}

Le document A, bien classé par les deux systèmes sans être premier partout,
devance le document C, premier pour un seul système. La fusion récompense donc
l'accord entre systèmes plutôt que l'excellence isolée.

\subsubsection{Variantes de la fusion}

\begin{btptable}{Familles de méthodes de fusion de classements}
\label{tab:familles-fusion}
\begin{tabularx}{\linewidth}{@{}L{3.1cm} Y Y@{}}
\toprule
\rowcolor{btppale}
\btpth{Méthode} & \btpth{Avantages} & \btpth{Inconvénients} \\
\midrule
Combinaison de scores & Théoriquement optimale si les scores sont calibrés. & Exige une normalisation dont le choix est arbitraire. \\
Comptage de Borda & Simple, fondée sur l'agrégation de préférences. & Perd l'information de score et l'écart entre rangs. \\
Fusion par rang réciproque & Sans normalisation, robuste, performante. & Un paramètre $k$ à régler. \\
Apprentissage du classement & Optimale en présence de données annotées. & Exige un corpus de jugements de pertinence. \\
\bottomrule
\end{tabularx}
\end{btptable}

\subsubsection{Le ré-ordonnancement}

Au-delà de la fusion, une étape supplémentaire affine le classement. Un
\btpterm{modèle croisé} (\anglais{cross-encoder}) encode conjointement la
requête et le document candidat, ce qui lui permet de modéliser leurs
interactions terme à terme. Il produit une estimation de pertinence bien plus
fine qu'une similarité entre vecteurs calculés indépendamment.

Ce gain a un coût~: là où une recherche vectorielle compare un vecteur de
requête à un index, le modèle croisé exige une inférence par couple. Il n'est
donc applicable qu'à un nombre restreint de candidats, ce qui impose
l'architecture en entonnoir de la figure~\ref{fig:pipeline-hybride}.

\begin{figure}[htbp]
  \centering
  \resizebox{\linewidth}{!}{% Pipeline de recherche hybride : BM25 + dense, fusion RRF, ré-ordonnancement
\begin{tikzpicture}[node distance=5mm and 10mm]

  \node[btpboxstrong, text width=2.4cm] (q) {Requête};

  \node[btpbox, text width=3.1cm, above right=2mm and 12mm of q] (bm25)
    {Recherche lexicale\\[-1pt]\btplabel{BM25 — termes exacts}};
  \node[btpbox, text width=3.1cm, below right=2mm and 12mm of q] (dense)
    {Recherche sémantique\\[-1pt]\btplabel{vecteurs denses — concepts}};

  \node[btpboxplain, text width=1.5cm, right=6mm of bm25] (t1) {Top 50};
  \node[btpboxplain, text width=1.5cm, right=6mm of dense] (t2) {Top 50};

  \node[btpboxalt, text width=2.4cm, right=8mm of {$(t1)!0.5!(t2)$}] (rrf)
    {Fusion\\\textbf{RRF}};
  \node[btpbox, text width=2.6cm, right=9mm of rrf] (rerank)
    {Ré-ordonnancement\\[-1pt]\btplabel{modèle croisé}};
  \node[btpboxstrong, text width=2.1cm, right=9mm of rerank] (out)
    {Top 8\\contexte};

  \draw[btparrow] (q) -- (bm25);
  \draw[btparrow] (q) -- (dense);
  \draw[btparrow] (bm25) -- (t1);
  \draw[btparrow] (dense) -- (t2);
  \draw[btparrow] (t1.east) -- (rrf.north west);
  \draw[btparrow] (t2.east) -- (rrf.south west);
  \draw[btparrow] (rrf) -- (rerank);
  \draw[btparrow] (rerank) -- (out);

  \node[btplabel, below=8mm of rrf, text width=3.6cm]
    {Aucune normalisation de score requise};
  \node[btplabel, below=8mm of rerank, text width=3.4cm]
    {Pertinence évaluée par couple (requête, extrait)};

  \begin{scope}[on background layer]
    \node[btpgroup, fit=(bm25)(dense)(t1)(t2),
          label={[btplabel]above:Récupération parallèle}] {};
  \end{scope}

\end{tikzpicture}
}
  \caption{Pipeline de recherche hybride en entonnoir}
  \label{fig:pipeline-hybride}
\end{figure}

\subsection{Application aux appels d'offres}

Les appels d'offres présentent des caractéristiques qui justifient une approche
hybride, résumées au tableau~\ref{tab:ri-ao}.

\begin{btptable}{Caractéristiques des appels d'offres et implications pour la recherche}
\label{tab:ri-ao}
\begin{tabularx}{\linewidth}{@{}L{3.9cm} Y@{}}
\toprule
\rowcolor{btppale}
\btpth{Caractéristique} & \btpth{Implication} \\
\midrule
Références normatives et identifiants & Recherche lexicale indispensable~: un numéro d'article ne se paraphrase pas. \\
Concepts métier exprimés diversement & Recherche sémantique indispensable~: la même exigence se formule de multiples façons. \\
Variantes de désignation des entités & Résolution d'entités nécessaire (section~\ref{sec:ch1-entites}). \\
Exigence de citation des sources & Précision de la récupération critique~: une citation erronée est plus dommageable qu'une abstention. \\
Structure hiérarchique des pièces & Le granule de recherche doit respecter la structure (section~\ref{sec:ch1-documents}). \\
\bottomrule
\end{tabularx}
\end{btptable}

Cette complémentarité justifie l'adoption d'un pipeline hybride combinant
recherche lexicale, recherche sémantique, fusion par rang réciproque et
ré-ordonnancement, dont la mise en \oe uvre est détaillée au
chapitre~\ref{chap:methodologie}.
